% ---- pg-588 （书页 576） ----

10. 监控模型的运行. 考察结果是否可靠, 成本是否合算?

在实践中, 经常有必要重复某些步骤若干次, 直至得到满意的结果为止. 数据挖掘软件套件, 如 Enterprise Miner 和 Clementine 等, 通常按照上述的步骤序列设计, 并且可以在屏幕上以流程图的方式图示结果.

数据挖掘需要丰富的工具和算法. 分析者不仅要能熟练使用这些工具和算法, 而且还要对研究课题的内涵有扎实的知识(或者与对研究课题的内涵有扎实知识的人合作), 如此才能得到可接受的结果. 任何数据挖掘结果一旦得到便成为一个不断进行的动态过程. 要采集新数据并对之加工, 模型要更新或开发新模型, 并且根据新的经验加以调整. 糟糕的数据挖掘工作的成本很高, 因此必须进行谨慎的模型建构和评估.

\mvsub{模型评估}

在数据挖掘的建模阶段, 可能同时检验若干模型. 在下面的案例中, 我们简要讨论对一组公开的数据——德国信用数据——用各种方法解决信用评分问题(确定信用风险大小)的结果. 这些方法是: 逻辑斯蒂回归, 决策树方法, 以及神经网络. 固然数据挖掘者能控制制约单个模型开发的模型输入和某些参数, 但是在大多数数据挖掘的努力中很少有正式的统计推断. 模型的评估与比较通常是由相关领域的专家来作, 使用的工具基本是描述性的, 例如混淆矩阵、综合效益或损失数量、提升图、阈值控制图以及其他方法, 其中大部分是图方法.

将巨大的原始数据集划分为训练集、验证集和检验集, 就可以在评估潜在模型时使用没有用于建模的数据. 这样, 训练数据集用于建模. 模型用验证(提留)数据集评估. 如果在验证阶段模型表现不够满意, 则暂时保留. 在训练和验证两个阶段反复几次直到得到满意的模型为止. 此时, 经过训练和验证的模型要用检验数据来评估. 检验数据通常用于建模过程的末尾, 以保证对模型的性能有无偏的评估. 有时检验数据步被省略, 此时最终的评估基于验证样本或基于交叉验证.

一种重要的评估估计是\mvbold{提升图}(lift chart). 提升图可以有不同的构造方法, 但是都指示所选择的程序(模型)相对于基准活动的改进. 基准活动通常表示某种先验的信念或者随机指定. 提升图对比较不同的模型特别有用.

提升(lift)定义为
\[
\mathrm{Lift}=P(\text{结果}\,|\,\text{条件})/P(\text{结果})
\]
若结果独立于条件, 则 Lift$=1$. 当 Lift$>1$, 则说明条件(通常是模型或算法)会使结果发生的概率增大. 因此条件是有用的并有潜在的效益. 不同的条件可以用各自的提升图来比较.

\noindent\textbf{例 12.23}\quad (小的数据挖掘练习)

一组可以公开得到的数据集称为德国信用数据.\textsuperscript{①}这组数据有 $1\,000$ 个信用申请人的 20 个变量的数据, 并且记录了每个申请人最终的信用等级(“好”或“坏”). (数据挖掘的)目标是开发一种信用评分法则用于决定任一新的申请人的信用风险是低还是高, 依据是 20 个解释变量的值. 20 个解释变量包括 CHECKING(支票账号状况), DURATION(以月为单位的信用期限), HISTORY(信用历史), AMOUNT(信用额), AGE(年龄), OTHERS(其他分期付款债务), INSTALLP(分期付款占储蓄收入的百分比), 等等. 本质上, 我们要开发一个关于若干个

\vspace{6pt}
\noindent\rule{\linewidth}{0.4pt}
\par\vspace{2pt}
{\zihao{6}\noindent ①\quad 在写作本补充材料时, 德国信用数据可从 SAS 的 Enterprise Miner 软件伴随的样本数据文件中得到. 许多公开的数据集可从网页 www.kdnuggets.com 上下载.\par}
